\section{Comparación de resultados}
\label{sec:comparacion}

\subsection{Evaluación con los mejores hiperparámetros}
\label{sec:evaluacion}

Los 10 puntos iniciales, generados con \texttt{random.Random(20260920)} y
distintos de los 6 puntos de ajuste, son los mismos para los tres
algoritmos y las tres funciones. Cada corrida usa los mejores
hiperparámetros de las Tablas~\ref{tab:gd_optuna}, \ref{tab:rms_optuna} y
\ref{tab:pso_optuna} y un máximo de 50 iteraciones. En las tablas
siguientes, ``Iter.'' es la primera iteración que cumple
$|f(x_t)-f^*|\leq10^{-3}$, o ``No'' si la corrida no convergió, y ``$f$''
es el valor en esa iteración o, si no convergió, al terminar las 50
iteraciones. Ninguna corrida divergió.

\begin{table}[H]
\centering
\small
\begin{tabular}{rlrrrrrr}
\toprule
 & & \multicolumn{2}{c}{GD} & \multicolumn{2}{c}{RMSProp} & \multicolumn{2}{c}{PSO} \\
\cmidrule(lr){3-4}\cmidrule(lr){5-6}\cmidrule(lr){7-8}
Corrida & Punto inicial & Iter. & $f$ & Iter. & $f$ & Iter. & $f$ \\
\midrule
1 & $(1.7280,\ 4.1962)$ & 2 & 2.71e-05 & 6 & 5.69e-04 & 32 & 3.58e-04 \\
2 & $(-7.4455,\ 8.0190)$ & 2 & 1.57e-04 & 3 & 3.86e-05 & 37 & 3.78e-04 \\
3 & $(-1.8004,\ -5.1750)$ & 2 & 3.94e-05 & 5 & 6.80e-04 & 34 & 9.40e-04 \\
4 & $(-1.2368,\ -9.5770)$ & 2 & 1.23e-04 & 5 & 1.26e-04 & No & 0.0028 \\
5 & $(4.4551,\ 0.2191)$ & 2 & 2.61e-05 & 6 & 3.67e-04 & No & 0.0021 \\
6 & $(-0.5907,\ 9.2344)$ & 2 & 1.13e-04 & 4 & 5.04e-04 & No & 0.0014 \\
7 & $(7.3363,\ -7.5118)$ & 2 & 1.45e-04 & 2 & 4.43e-05 & No & 0.0052 \\
8 & $(-6.2477,\ -1.4124)$ & 2 & 5.39e-05 & 5 & 7.20e-05 & 49 & 4.84e-04 \\
9 & $(-6.2994,\ -0.1350)$ & 2 & 5.22e-05 & 3 & 4.63e-04 & No & 0.0026 \\
10 & $(-0.4445,\ 7.9754)$ & 2 & 8.38e-05 & 3 & 8.03e-04 & No & 0.0013 \\
\bottomrule
\end{tabular}
\caption{Resultados de las 10 corridas sobre $f_0$.}
\label{tab:eval_f0}
\end{table}

\begin{table}[H]
\centering
\small
\begin{tabular}{rlrrrrrr}
\toprule
 & & \multicolumn{2}{c}{GD} & \multicolumn{2}{c}{RMSProp} & \multicolumn{2}{c}{PSO} \\
\cmidrule(lr){3-4}\cmidrule(lr){5-6}\cmidrule(lr){7-8}
Corrida & Punto inicial & Iter. & $f$ & Iter. & $f$ & Iter. & $f$ \\
\midrule
1 & $(1.7280,\ 4.1962)$ & No & 2.5883 & No & 9.3533 & No & 0.3863 \\
2 & $(-7.4455,\ 8.0190)$ & No & 12.8191 & No & 15.5491 & No & 0.0199 \\
3 & $(-1.8004,\ -5.1750)$ & No & 4.6796 & No & 10.6451 & No & 0.0925 \\
4 & $(-1.2368,\ -9.5770)$ & No & 12.7398 & No & 15.1676 & No & 0.0625 \\
5 & $(4.4551,\ 0.2191)$ & No & 3.9015 & No & 8.6165 & No & 0.4717 \\
6 & $(-0.5907,\ 9.2344)$ & No & 4.0646 & No & 14.4371 & No & 0.1134 \\
7 & $(7.3363,\ -7.5118)$ & No & 4.6046 & No & 15.0637 & No & 0.0620 \\
8 & $(-6.2477,\ -1.4124)$ & No & 7.2583 & No & 11.5254 & No & 0.1122 \\
9 & $(-6.2994,\ -0.1350)$ & No & 4.5613 & No & 11.4256 & No & 0.1062 \\
10 & $(-0.4445,\ 7.9754)$ & No & 3.9212 & No & 12.5578 & No & 0.0701 \\
\bottomrule
\end{tabular}
\caption{Resultados de las 10 corridas sobre $f_1$ (Ackley).}
\label{tab:eval_f1}
\end{table}

\begin{table}[H]
\centering
\small
\begin{tabular}{rlrrrrrr}
\toprule
 & & \multicolumn{2}{c}{GD} & \multicolumn{2}{c}{RMSProp} & \multicolumn{2}{c}{PSO} \\
\cmidrule(lr){3-4}\cmidrule(lr){5-6}\cmidrule(lr){7-8}
Corrida & Punto inicial & Iter. & $f$ & Iter. & $f$ & Iter. & $f$ \\
\midrule
1 & $(1.7280,\ 4.1962)$ & 33 & 8.51e-04 & No & 1.9068 & No & 0.1448 \\
2 & $(-7.4455,\ 8.0190)$ & 37 & 7.48e-04 & 24 & 6.73e-07 & No & 0.2348 \\
3 & $(-1.8004,\ -5.1750)$ & 14 & 4.33e-04 & 12 & 6.70e-05 & No & 0.0471 \\
4 & $(-1.2368,\ -9.5770)$ & 16 & 4.52e-04 & No & 3.8037 & No & 0.1478 \\
5 & $(4.4551,\ 0.2191)$ & No & 6.7753 & 7 & 8.29e-04 & No & 0.5138 \\
6 & $(-0.5907,\ 9.2344)$ & 19 & 7.20e-04 & No & 2.0763 & No & 0.2357 \\
7 & $(7.3363,\ -7.5118)$ & 33 & 8.62e-04 & No & 2.0284 & No & 0.1239 \\
8 & $(-6.2477,\ -1.4124)$ & 16 & 6.45e-04 & 14 & 3.61e-05 & No & 0.3294 \\
9 & $(-6.2994,\ -0.1350)$ & 26 & 5.45e-04 & 14 & 7.69e-04 & No & 0.1800 \\
10 & $(-0.4445,\ 7.9754)$ & 20 & 7.17e-04 & No & 2.2088 & No & 0.6862 \\
\bottomrule
\end{tabular}
\caption{Resultados de las 10 corridas sobre $f_2$ (Himmelblau).}
\label{tab:eval_f2}
\end{table}

\begin{table}[H]
\centering
\begin{tabular}{llcrr}
\toprule
Algoritmo & Función & Convergencias & Promedio $f$ minimizada & Iter.\ promedio \\
\midrule
GD & $f_0$ & 10/10 & 8.20e-05 & 2.00 \\
GD & $f_1$ & 0/10 & 6.1138 & --- \\
GD & $f_2$ & 9/10 & 0.6781 & 23.78 \\
RMSProp & $f_0$ & 10/10 & 3.67e-04 & 4.20 \\
RMSProp & $f_1$ & 0/10 & 12.4341 & --- \\
RMSProp & $f_2$ & 5/10 & 1.2026 & 14.20 \\
PSO & $f_0$ & 4/10 & 0.0018 & 38.00 \\
PSO & $f_1$ & 0/10 & 0.1497 & --- \\
PSO & $f_2$ & 0/10 & 0.2643 & --- \\
\bottomrule
\end{tabular}
\caption{Resumen de las 10 corridas. El promedio de $f$ usa las 10
corridas; el promedio de iteraciones usa solo las corridas convergentes.}
\label{tab:resumen_eval}
\end{table}

\subsection{Descenso del gradiente frente a RMSProp}

\begin{itemize}
    \item \textbf{$f_0$.} Ambos convergen en las 10 corridas. GD es más
    rápido (2 iteraciones frente a $4.20$) porque en una cuadrática
    isotrópica un único $\alpha$ cercano a $0.5$ es casi el paso óptimo; la
    normalización de RMSProp no aporta nada en ese caso.
    \item \textbf{$f_1$.} Ninguno converge. GD obtiene un promedio menor
    ($6.11$ frente a $12.43$) porque su $\alpha$ calibrado es relativamente grande y le
    permite saltar entre cuencas, mientras que RMSProp, con paso pequeño,
    queda en el mínimo local cercano a su inicio. La mejora de GD no es
    un descenso controlado: su curva de aprendizaje es no monótona
    (Figura~\ref{fig:gd_mejores}).
    \item \textbf{$f_2$.} GD converge en más corridas (9/10 frente a 5/10),
    pero RMSProp converge en menos iteraciones cuando lo hace ($14.20$
    frente a $23.78$) y su corrida más rápida necesitó 7 iteraciones frente a
    14 de GD. RMSProp tolera pasos 60 veces mayores ($0.341559$ frente a
    $0.005492$) sin divergir, y en la corrida 5 evitó el punto silla donde
    GD quedó detenido (Tabla~\ref{tab:silla}). Su debilidad es la
    precisión final: con $\gamma=0.611311$ el paso normalizado no decrece
    cerca del mínimo y 5 corridas terminan oscilando con $f$ entre $1.9$ y
    $3.8$.
\end{itemize}

En resumen, con el presupuesto de 25 iteraciones de calibración y 50 de
evaluación, GD fue más confiable en $f_0$ y $f_2$ y RMSProp fue más rápido
en las corridas de $f_2$ que sí convergieron. RMSProp es menos sensible a la
escala del gradiente, pero su paso casi constante dificulta refinar la
solución si no se reduce $\alpha$ con el tiempo.

\subsection{Comparación de los tres algoritmos}

\begin{itemize}
    \item \textbf{Convergencia a la tolerancia.} Solo GD y RMSProp alcanzan
    $10^{-3}$ de forma sistemática en $f_0$, y GD en 9 de 10 corridas de
    $f_2$. PSO converge en 4 corridas de $f_0$ y en ninguna de $f_1$ ni
    $f_2$.
    \item \textbf{Calidad promedio.} En las funciones no convexas PSO obtiene
    el menor promedio de $f$: $0.1497$ en $f_1$ frente a $6.11$ (GD) y
    $12.43$ (RMSProp), y $0.2643$ en $f_2$ frente a $0.678$ y $1.203$. Todas
    sus corridas en $f_1$ terminan con $f<0.48$, por debajo del mínimo
    local más bajo de Ackley ($2.5799$), es decir, en la cuenca del mínimo
    global.
    \item \textbf{Velocidad.} Cuando convergen, GD y RMSProp necesitan de 2
    a 37 iteraciones; PSO necesita de 32 a 49 en $f_0$.
    \item \textbf{Costo por iteración.} GD y RMSProp evalúan $f$ y su
    gradiente en un punto; PSO evalúa $f$ en 36 a 43 partículas, aunque de
    forma vectorizada.
\end{itemize}

\subsection{Ventajas y desventajas}

\paragraph{Descenso del gradiente.}
\emph{Ventajas:} es el método más simple, con un solo hiperparámetro, y en
funciones suaves y bien condicionadas converge rápido cerca de un mínimo
\cite{nocedal2006numerical}; en nuestros resultados fue el más rápido y
confiable en $f_0$ y el más confiable en $f_2$. \emph{Desventajas:} es un
método local que converge al mínimo de la cuenca donde inicia; es muy
sensible a $\alpha$, que debe ajustarse a la escala del gradiente (en $f_2$,
9 de 25 ensayos de Optuna divergieron); y se desacelera cerca de puntos
silla y en mesetas \cite{goodfellow2016deep,dauphin2014identifying}, como
en la corrida 5 sobre $f_2$.

\paragraph{RMSProp.}
\emph{Ventajas:} adapta el tamaño de paso por coordenada con un promedio
móvil de gradientes al cuadrado, lo que lo hace menos sensible a la escala
del gradiente y a problemas mal condicionados, y le ayuda a atravesar
mesetas y puntos silla \cite{tieleman2012rmsprop,ruder2016overview}; en
$f_2$ usó pasos mucho mayores sin divergir y convergió más rápido cuando lo
hizo. \emph{Desventajas:} agrega hiperparámetros ($\gamma$, $\epsilon$);
sigue siendo un método local que no escapa de mínimos locales
\cite{goodfellow2016deep}, como muestran sus resultados en Ackley; y, con
$\alpha$ fijo, el paso normalizado no se reduce cerca del mínimo, lo que
produce oscilaciones.

\paragraph{Enjambre de partículas.}
\emph{Ventajas:} no requiere gradientes, por lo que sirve para funciones no
diferenciables como Ackley en el origen; su búsqueda poblacional explora
el dominio y es más robusta ante mínimos locales
\cite{kennedy1995particle}, como muestra su promedio en $f_1$ y $f_2$; y se
vectoriza de forma natural. \emph{Desventajas:} evalúa $f$ muchas veces por
iteración; su convergencia fina es lenta y depende del ajuste de $c_1$,
$c_2$, $n$ y, en variantes con inercia, $w$
\cite{shi1998modified,clerc2002particle}; no garantiza convergencia a un
mínimo; y, en la forma sin inercia implementada, las partículas siguen
oscilando, por lo que no alcanzó la tolerancia de $10^{-3}$ en $f_1$ ni en
$f_2$.
